% ============================================================================
% fig_diagrama_2_schur_hbcf.tex
% DIAGRAMA 2 — Desacoplo diakóptico HBCF y complemento de Schur en el SENI
%              insular: 142 barras -> 4 subdominios, fronteras Gamma, KPIs en
%              silicio. Informe hito SENI 142 barras (Rep. Dominicana).
% Dossier TPWRS-SENI-142-2026 · Autoría: Ing. J. A. Villanueva Rosario
% ----------------------------------------------------------------------------
% Compilación individual (protocolo solucion_error_latex.md):
%   nice -n 19 pdflatex -interaction=nonstopmode fig_diagrama_2_schur_hbcf.tex
%
% Cumple: regla 1 (arrows.meta + tip explícito, nunca > ni < crudos),
% regla 2 (align=center + text width/minimum height + rounded corners),
% regla 3 (paleta doctoral soberana), regla 4 (cero menciones prohibidas),
% regla 5 (bloque autocontenido). Cifras: únicamente las del informe hito.
% ============================================================================
\documentclass[border=12pt]{standalone}
\usepackage[T1]{fontenc}
\usepackage[utf8]{inputenc}
\usepackage[spanish,es-nodecimaldot,es-noshorthands]{babel}
\usepackage{amsmath,amssymb}
\usepackage{tikz}
\usetikzlibrary{arrows.meta,positioning,calc,fit,backgrounds,babel}

\begin{document}
% ==================== BLOQUE INSERTABLE EN IEEETRAN: INICIO ====================
% Requisitos del documento hospedador (p. ej. IEEEtran):
%   \usepackage{amsmath,amssymb}
%   \usepackage{tikz}
%   \usetikzlibrary{arrows.meta,positioning,calc,fit,backgrounds,babel}
\tikzset{>={Stealth[length=2.5mm,width=1.8mm]}}
\definecolor{seniNavy}{RGB}{26,54,93}   % Azul institucional
\definecolor{seniBlue}{RGB}{43,108,176} % Azul alta tensión
\definecolor{seniGold}{RGB}{217,194,130}% Dorado V-Core
\definecolor{seniRed}{RGB}{155,44,44}   % Rojo 345 kV / fronteras Gamma
\definecolor{seniGreen}{RGB}{39,103,73} % Verde renovable
\definecolor{seniBg}{RGB}{247,250,252}  % Fondo de caja
\begin{tikzpicture}[
  font=\small,
  region/.style={align=center, text width=46mm, minimum height=20mm,
                 rounded corners=3pt, line width=1pt, font=\footnotesize},
  gamma/.style={midway, fill=white, inner sep=1.2pt, font=\scriptsize\bfseries,
                text=seniRed},
  nota/.style={align=center, text width=42mm, minimum height=12mm,
               rounded corners=2pt, draw=seniRed, line width=0.7pt, dashed,
               fill=white, font=\tiny, text=seniRed},
  paso/.style={align=center, text width=44mm, minimum height=16mm,
               rounded corners=2pt, draw=seniNavy, line width=1pt,
               fill=seniBg, font=\footnotesize},
  kpi/.style={align=center, text width=40mm, minimum height=22mm,
              rounded corners=3pt, draw=seniGold!60!black, line width=1.2pt,
              fill=seniNavy, text=seniGold, font=\footnotesize},
  titulo/.style={align=center, text width=176mm, minimum height=8mm,
                 rounded corners=2pt, draw=seniNavy, fill=seniNavy,
                 text=white, font=\small\bfseries},
  rotulo/.style={align=center, font=\scriptsize, text=seniNavy}
]

% ---- Cintillo de título ----
\node[titulo] (tit) at (9.5, 12.3)
  {Desacoplo diakóptico HBCF del SENI insular: 142 barras $\to$ cuatro
   subdominios $\Omega$, fronteras $\Gamma$ y complemento de Schur};

% ---- Mapa topológico insular: 4 macro-regiones desacopladas ----
\node[region, draw=seniBlue, fill=seniBlue!8] (norte) at (3.2, 9.9)
  {$\Omega_{\text{Norte}}$ --- Cibao\\[0.6mm]
   {\scriptsize Canabacoa · Tavera · Monción\\ Falcondo · Pueblo Viejo}};
\node[region, draw=seniGreen, fill=seniGreen!8] (este) at (14.9, 9.3)
  {$\Omega_{\text{Este}}$\\[0.6mm]
   {\scriptsize San Pedro · Quisqueya · AES Andrés\\ Corredor Turístico Punta Cana}};
\node[region, draw=seniRed, fill=seniRed!8] (sur) at (3.2, 5.6)
  {$\Omega_{\text{Sur}}$\\[0.6mm]
   {\scriptsize Punta Catalina · Julio Sauri · Barahona\\
    Parques solares/eólicos · PHES El Valle}};
\node[region, draw=seniNavy, fill=seniNavy!8] (metro) at (11.4, 4.9)
  {$\Omega_{\text{Metro}}$ --- Gran Santo Domingo\\[0.6mm]
   {\scriptsize Palamara · Hainamosa · Herrera}\\
   {\footnotesize\bfseries \textgreater{}40\,\% de la carga}};

% ---- Nota de fronteras diakópticas (hueco central izquierdo) ----
\node[nota] (notag) at (3.2, 7.75)
  {$\Gamma$: enlaces troncales 230/138 kV cortados virtualmente
   (descomposición de Kron)};

% ---- Panel de KPIs en silicio ----
\node[kpi] (kpi) at (16.6, 5.4)
  {\textbf{KPIs en silicio}\\[0.8mm]
   Tiempo: \textbf{2.65 ms}/flujo\\[0.3mm]
   RAM: \textbf{473 MiB}\\[0.3mm]
   Horizonte: \textbf{87,360 horas}\\[0.3mm]
   {\scriptsize (10 años OC-SENI)}};

% ---- Fronteras Gamma entre regiones (cortes de Kron) ----
\draw[dashed, line width=1.4pt, seniRed] (norte) -- (metro)
  node[gamma]{$\Gamma$};
\draw[dashed, line width=1.4pt, seniRed] (norte) -- (este)
  node[gamma]{$\Gamma$};
\draw[dashed, line width=1.4pt, seniRed] (sur) -- (metro)
  node[gamma]{$\Gamma$};
\draw[dashed, line width=1.4pt, seniRed] (metro) -- (este)
  node[gamma]{$\Gamma$};

% ---- Flujo algorítmico diakóptico (banda inferior) ----
\node[paso] (p1) at (2.9, 1.5)
  {\textbf{1 · Factorización paralela}\\[0.4mm]
   Jacobianas ralas locales\\ $\mathbf{J}_1,\ \mathbf{J}_2,\ \mathbf{J}_3,\ \mathbf{J}_4$};
\node[paso] (p2) at (9.5, 1.5)
  {\textbf{2 · Complemento de Schur en $\Gamma$}\\[0.4mm]
   $\mathbf{S}=\mathbf{D}-\mathbf{C}\,\mathbf{A}^{-1}\mathbf{B}$};
\node[paso] (p3) at (16.0, 1.5)
  {\textbf{3 · Reconstrucción nodal}\\[0.4mm]
   tensiones completas\\ acople fasorial nulo\\
   $\delta=2.788\times10^{-14}$};

\draw[->, line width=1.6pt, seniNavy] (p1) -- (p2);
\draw[->, line width=1.6pt, seniNavy] (p2) -- (p3);

% ---- Las fronteras Gamma alimentan la reducción de Schur ----
\draw[->, dashed, line width=1.1pt, seniRed] (metro.south) -- (p2.north)
  node[rotulo, pos=0.45, right=1mm]{fronteras $\Gamma$};

\end{tikzpicture}
% ===================== BLOQUE INSERTABLE EN IEEETRAN: FIN =====================
\end{document}
